% MCDS Capstone Plan Document
% Kyle Dotterrer
% 15 March, 2021

%% ==================================================================
%% BOILERPLATE
%% ==================================================================

\documentclass{vldb}

\newcommand{\paperTitle}{JIT Compilation for User-Defined Functions\\ 
    \large Capstone Plan Document\\
    \today}
\newcommand{\paperKeywords}{}
\newcommand{\paperAuthors}{}

\setlength{\paperheight}{11in}
\setlength{\paperwidth}{8.5in}

\usepackage[
            bookmarks=true,
            bookmarksopen=true,
            pdfhighlight=/I,
            pdfpagemode=UseOutlines,
            linkcolor=blue,
            pdfborder={ 0 0 0 },
            pageanchor=false]{hyperref}
\hypersetup{
    pdfauthor = {\paperAuthors},
    pdftitle = {\paperTitle},
    pdfkeywords = {\paperKeywords},
    pdfborder={ 0 0 0 }
}

\usepackage{amsmath}
\usepackage{amssymb}
\usepackage{times}
\usepackage{boxedminipage}
\usepackage{xspace}
\usepackage{array}
\usepackage{epsfig}
\usepackage{calc}
\usepackage{multirow}
\usepackage{rotating}
\usepackage{enumitem}
\usepackage[usenames,dvipsnames]{color}
\usepackage{tabularx}
\usepackage{tikz}
\usepackage{wasysym }
\usepackage[hyphenbreaks]{breakurl}
\usepackage{xcolor}
\usepackage[font={small}]{caption}
\usepackage{alltt}
\usepackage{setspace}
\usepackage{wrapfig}
\usepackage{etoolbox}

% Font selection
\usepackage[final]{microtype}
\usepackage[scaled]{inconsolata}
\usepackage[T1]{fontenc}
\usepackage{graphicx}
\usepackage{soul}
\usepackage{pifont}

% Other misc. things
\usepackage{float}
\usepackage{bookmark}
\usepackage{amsmath,amsfonts,amssymb}
\usepackage{latexsym}
\usepackage{multirow}
\usepackage{emptypage}
\usepackage{booktabs}
\usepackage{enumerate}
\usepackage{braket}
\usepackage{balance}

%% ==================================================================
%% MACROS
%% ==================================================================

\newcommand{\subparagraph}{}

\newcommand{\sysname}{NoisePage\xspace}

% Framework Components
\newcommand{\dbTranslator}{Translator\xspace}
\newcommand{\dbCompiler}{Compiler\xspace}
\newcommand{\dbPlanner}{planner\xspace}
\newcommand{\dbInterpreter}{interpreter\xspace}
\newcommand{\dbAnalyzer}{analyzer\xspace}
\newcommand{\dbPCQ}{PCQ\xspace}
\newcommand{\dbAQP}{AQP\xspace}
\newcommand{\dbHook}{hook\xspace}
\newcommand{\dbHooks}{hooks\xspace}

\newcommand{\dbSQL}[1]{\texttt{\textbf{#1}}\xspace}
\newcommand{\dbTable}[1]{\texttt{\MakeUppercase{#1}}\xspace}
\newcommand{\dbColumn}[1]{\texttt{col#1}\xspace}
\newcommand{\dbRawColumn}[1]{\texttt{#1}\xspace}
\newcommand{\dbConfig}[1]{\textsf{#1}\xspace}

\newcommand{\dbmsOracle}{Oracle\xspace}
\newcommand{\dbmsDBTwo}{DB2\xspace}
\newcommand{\dbmsTeradata}{Teradata\xspace}
\newcommand{\dbmsHANA}{HANA\xspace}
\newcommand{\dbmsMSSQL}{SQL Server\xspace}
\newcommand{\dbmsMySQL}{MySQL\xspace}
\newcommand{\dbmsPostgres}{PostgreSQL\xspace}
\newcommand{\dbmsMemSQL}{MemSQL\xspace}
\newcommand{\dbmsHyper}{HyPer\xspace}
\newcommand{\dbmsVector}{Vector\xspace}

\newcommand{\tplType}[1]{\texttt{#1}\xspace}
\newcommand{\tpl}{\texttt{TPL}\xspace}
\newcommand{\tbc}{\texttt{TBC}\xspace}
\newcommand{\dbCode}[1]{{\sffamily\small \textbf{#1}}\xspace}

\newcommand{\tidlist}{\texttt{TupleIdList}\xspace}
\newcommand{\tidlists}{\texttt{TupleIdLists}\xspace}

%% ==================================================================
%% PAPER CONTENT
%% ==================================================================

\begin{document}

\title{\paperTitle}

\numberofauthors{2}
\author{
    \alignauthor Kyle Dotterrer, Andrew Pavlo\\
    \affaddr{Carnegie Mellon University}\\
    \email 
        \{\href{mailto:kdottere@cs.cmu.edu}{kdottere},
        \href{mailto:pavlo@cs.cmu.edu}{pavlo}\}@cs.cmu.edu\\
}

% Remove the space for the copyright box for VLDB
\makeatletter
\def\@copyrightspace{\relax}
\makeatother

\maketitle

%% ==================================================================
%% INTRODUCTION
%% ==================================================================

\section{Introduction}

User-defined functions (UDFs) are a powerful tool for encapsulating complex application logic that may subsequently be embedded into SQL queries. Although they offer a practical complement to SQL, their use in production database environments is often limited because of their associated performance costs \cite{froid, apfel}. In this project, we aim to address the performance limitations of UDFs in the context of main-memory database management systems (DBMSs). We approach this issue by implementing a novel UDF execution technique via JIT-compilation and evaluating this implementation's performance against existing solutions. This document describes our high-level plan for meeting our project's goals.

Our project is naturally partitioned into three distinct components or phases. The first two of these phases consist of system design and implementation tasks, while the third consists of experimental evaluation. We provide additional details regarding each of these tasks in the sections that follow.

%% ==================================================================
%% TASKS
%% ==================================================================

\section{Task: UDF JIT-Compilation}

The first task in our project is the design and implementation of a subsystem for UDF execution via JIT-compilation in NoisePage \cite{noisepage}.

\textbf{Task Description}. In this task, we implement a novel execution strategy for user-defined functions in the research database management system NoisePage. This is a system design and implementation task.

\textbf{Start Date and End Date}. This task is ongoing. Our goal is to complete this task by 1 May 2021.

\textbf{Required Resources}. There are no resources required for this task other than the development effort involved.

\textbf{Task Dependencies}. This task is not dependent on the completion of other tasks.

\section{Task: UDF Inlining}

The second task in our project is the design and implementation of a subsystem for UDF execution via UDF inlining in NoisePage \cite{noisepage}.

\textbf{Task Description}. In this task, we implement support for inlining of user-defined functions in the research database management system NoisePage. This is a system design and implementation task.

\textbf{Start Date and End Date}. This task will begin once the UDF JIT-compilation implementation task is complete. Our goal is to begin this task by 1 May 2021. Our goal is to complete this task by 1 June 2021.

\textbf{Required Resources}. There are no resources required for this task other than the development effort involved.

\textbf{Task Dependencies}. This task is not dependent on the completion of other tasks.

\section{Task: Experimental Evaluation}

The final task required to meet our project's objectives is an experimental evaluation of the features implemented in the prior two tasks.

\textbf{Task Description}. In this task, we perform an evaluation of performance and resource consumption characteristics of various execution techniques for user-defined functions. This is an experimentation task.

\textbf{Start Date and End Date}. This task will begin once both the UDF JIT-compilation and UDF inlining implementation tasks are complete. Our goal is to begin this task by 1 June 2021. Our goal is to complete this task by 1 August 2021.

\textbf{Required Resources}. This task requires access to workload characteristic data and the associated queries in order to construct benchmarks for evaluation.

\textbf{Task Dependencies}. This task is dependent on both the UDF JIT-compilation task as well as the UDF inlining task.

%% ==================================================================
%% MILESTONE PLAN
%% ==================================================================

\section{Milestone Plan}

The project milestones that we have identified are enumerated in the table below. The dates listed are estimates at best; it is likely that the project timeline will shift as we encounter advances and setbacks. It is possible that we underestimate the difficulty of the system design and implementation tasks identified above.

\subsection{UDF JIT Compilation}

\begin{center}
\begin{tabular}{ |c|l| } 
    \hline
    \textbf{Date} & \textbf{Milestone} \\ 
    \hline
    1 March 2021 & Begin System Design \\
    \hline
    21 March 2021 & All tests for feature designed \\
    \hline
    28 March 2021 & \textit{Parser} support implemented; tests passing \\
    \hline
    4 April 2021 & \textit{Binder} support implemented; tests passing \\
    \hline
    11 April 2021 & \textit{Execution Engine} codegen implemented; tests passing \\
    \hline
    18 April 2021 & \textit{Execution Engine} LLVM implemented; tests passing \\
    \hline
    25 April 2021 & Integration tests implemented and passing \\
    \hline
    1 May 2021 & Feature Complete \\
    \hline
\end{tabular}
\end{center}

\subsection{UDF Inlining}

\begin{center}
\begin{tabular}{ |c|l| } 
    \hline
    \textbf{Date} & \textbf{Milestone} \\ 
    \hline
    1 May 2021 & Begin system design \\
    \hline
    9 May 2021 & System design complete; integration plan \\
    \hline
    16 May 2021 & \textit{Binder} support implemented; tests passing \\
    \hline
    23 May 2021 & \textit{Optimizer} support implemented; tests passing \\
    \hline  
    1 June 2021 & Feature Complete \\
    \hline
\end{tabular}
\end{center}

\subsection{Experimental Evaluation}

\begin{center}
\begin{tabular}{ |c|l| } 
    \hline
    \textbf{Date} & \textbf{Milestone} \\ 
    \hline
    1 June 2021 & Begin Experimental Setup \\
    \hline
    6 June 2021 & Complete experimental design for all cases \\
    \hline
    13 June 2021 & All workload data collected \\
    \hline
    20 June 2021 & All targeted benchmarks implemented \\
    \hline
    27 June 2021 & All production benchmarks implemented \\
    \hline
    4 July 2021 & Targeted benchmarks complete \\
    \hline
    11 July 2021 & Production benchmarks complete \\
    \hline
    18 July 2021 & Benchmark results organized into tables \\
    \hline
    25 July 2021 & Initial evaluation of results complete \\
    \hline
    1 August 2021 & Experimentation Complete \\
    \hline
\end{tabular}
\end{center}

%% ==================================================================
%% BIBLIOGRAPHY
%% ==================================================================

% NOTE: Had to add this mysterious command to compile...
\newcommand{\newblock}{}

\newpage
\bibliographystyle{abbrv}
\bibliography{plan}

\end{document}